\ExerciseSection

¶ 1) 对复数的每个有限序列 \(s = (z_k)_{k \in \mathbf{I}}\) ，令

\[
\sum_ {k \in \mathbf {I}} | z _ {k} | = \rho_ {s} \cdot \sup _ {\mathbf {J} \subset \mathbf {I}} \left| \sum_ {k \in \mathbf {J}} z _ {k} \right|
\]

（参看第七章，§ 3，第 1 号，命题 2）。

\begin{enumerate}
\def\labelenumi{\alph{enumi})}
\tightlist
\item
  * 若 \(z_{k} = r_{k} (\cos \varphi_{k} + i \sin \varphi_{k})\) ，其中 \(r_{k} = |z_{k}|\) ，而 \(\varphi_{k}\) 是 \(z_{k}\) 的辐角的主弧度测度，令
\end{enumerate}

\[
f (\varphi) = \sum_ {k \in \mathbf {I}} r _ {k} (\cos (\varphi - \varphi_ {k})) ^ {+}
\]

对每个 \(\varphi \in [0, 2\pi[\) 成立。证明 \(f(\varphi)\) 在区间 \([0, 2\pi[\) 中的上确界等于 \(\sup_{J \subset I} \left| \sum_{k \in J} z_k \right|\) 。考察积分 \(\int_0^{2\pi}f(\varphi)d\varphi\) ，推出：对一切有限序列 \(s\) 有 \(\rho_{s}\leqslant \pi\) ，且不存在有限序列 \((z_k) = s\) 使得 \(\rho_{s} = \pi\)

\begin{enumerate}
\def\labelenumi{\alph{enumi})}
\setcounter{enumi}{1}
\tightlist
\item
  证明：对每个 \(\varepsilon > 0\) ，存在有限序列 \((z_k) = s\) 使得 \(\rho_s > \pi - \varepsilon\) （取 \(z_k\) 为一个次数足够高的二项方程的根）。*
\end{enumerate}

\begin{enumerate}
\def\labelenumi{\arabic{enumi})}
\setcounter{enumi}{1}
\tightlist
\item
  设 \((z_{n})\) 是复数的无穷序列 \(z_{n} = x_{n} + iy_{n}\) ，使得 \(x_{n} \geqslant 0\) 对一切 \(n\) 成立。证明：若以 \(z_{n}\) 与 \(z_{n}^{2}\) 为通项的级数都收敛，则以 \(z_{n}^{2}\) 为通项的级数绝对收敛。给出一个例子，说明当把条件 \(x_{n} \geqslant 0\) （它也可写成 \(-\pi / 2 \leqslant \operatorname{Am}(z_{n}) \leqslant \pi / 2\) ）换成
\end{enumerate}

\[
- (\mathrm{I} - \varepsilon) \frac {\pi}{2} \leqslant \operatorname{Am} (z _ {n}) \leqslant (\mathrm{I} + \varepsilon) \frac {\pi}{2}
\]

时，无论实数 \(\varepsilon > 0\) 多么小，该结果都可能不成立。（这里 Am \((z)\) 表示 \(z_{n}\) 的辐角的属于区间 \([- \pi, +\pi]\) 的那个弧度测度。）

\begin{enumerate}
\def\labelenumi{\arabic{enumi})}
\setcounter{enumi}{2}
\item
  设 \((z_{t})_{t\in I}\) 是一族复数，使得 \(\prod_{t\in I}|z_t| = +\infty\) （相应地 \(\prod_{t\in I}|z_t| = 0\) ）。证明族 \((z_{t})\) 在空间 \(\tilde{\mathbf{C}}\) （ \(\S 4\) ，第 3 号）中可乘，且其乘积为 \(\infty\) （相应地为 0）。
\item
  证明以 \(1 + i/n\) 为通项因子的无穷乘积不收敛，但其因子的绝对值之乘积绝对收敛。
\item
  设 \((z_{n})\) 是非零复数的无穷序列，使得 \(\lim_{n\to \infty}z_n = 1\) 。证明：若 \(\sigma\) 是 \(\mathbf{N}\) 的一个置换，使得以 \(z_{\sigma (n)}\) 为通项因子的无穷乘积收敛，则诸乘积
\end{enumerate}

\[
\underset {n = 0} {\overset {\infty} {\mathbf {P}}} z _ {\sigma (n)}
\]

所成的集，当取遍所有这些置换时，可以是：(i) 单个点；或 (ii) 整个群 \(\mathbf{C}^*\) ；或 (iii) 以 \(\mathbf{o}\) 为端点的一条开射线；或 (iv) 以 \(\mathbf{o}\) 为圆心的一个圆周；或 (v) 一条“对数螺线”，即 \(\mathbf{R}\) 在映射 \(t \to a^{t - t_0} (\cos t + i \sin t)\) 下的像，其中 \(a\) 是 \(>0\) 且 \(\neq 1\) 的实数。（利用乘法群 \(\mathbf{C}^*\) 同构于加法群 \(\mathbf{R} \times \mathbf{T}\) 这一事实，按第七章 § 3 习题 2 的方式论证。）

